% Ch3 supplement -- elaborate worked-example study notes for Zumdahl
% 3.7-3.11. Complements ch03-notes (the block lessons): every method is
% carried through a full worked example at reading pace, then self-checked.
% Compounds deliberately differ from those used in ch03-notes.
\documentclass{apchem-worksheet}

\apcunitword{Chapter}
\apcunit{3}
\apcunittitle{Stoichiometry}
\apcdoctitle{Study Notes \textbullet\ \S3.7--3.11 in Worked Examples}
\apcsubtitle{Zumdahl \S3.7--3.11 \textbullet\ PDF pp.\ 125--153 \textbullet\ read with a pencil}

\begin{document}
\apctitleblock

\begin{apcnote}[What these five sections are, as one story]
Section 3.7 answers ``what \emph{is} this compound?'' --- turning measured
percentages or combustion data into a formula. Sections 3.8--3.9 give the
grammar for writing what compounds \emph{do}: the balanced equation.
Sections 3.10--3.11 then use that equation as a conversion factor: given an
amount of one substance, find the amount of any other --- including the
case where one reactant runs out first.

Every calculation in this span is the same three-step spine:
\[ \text{grams} \xrightarrow{\;\div M\;} \text{moles}
   \xrightarrow{\;\text{mole ratio}\;} \text{moles}
   \xrightarrow{\;\times M\;} \text{grams} \]
Only the middle step ever uses the chemical equation. Grams never talk to
grams directly.
\end{apcnote}

% =====================================================================
\section*{\sffamily \S3.7 \textbullet\ Determining the formula of a
compound}
\zsec{3.7}

A formula determination always runs the same four steps:

\begin{enumerate}[leftmargin=*,itemsep=0.2em]
  \item Assume a \SI{100}{\gram} sample, so percentages become grams.
  \item Convert each element's mass to moles.
  \item Divide every mole number by the smallest one.
  \item If the ratios are not whole, multiply \emph{all} of them by the
        same small integer. The result is the \term{empirical formula};
        the molecular formula is the empirical formula times
        $M / M_{\text{empirical}}$.
\end{enumerate}

\begin{workedexample}[Worked example 1: from percentages --- nicotine]
Nicotine is 74.03\% C, 8.70\% H, and 17.27\% N, with a molar mass of
\SI{162.23}{\gram\per\mole}. Find both formulas.

\textbf{Step 1 --- take \SI{100}{\gram}:} 74.03 g C, 8.70 g H, 17.27 g N.

\textbf{Step 2 --- to moles:}
\[ \ce{C}: \frac{74.03}{12.01} = 6.164 \qquad
   \ce{H}: \frac{8.70}{1.008} = 8.63 \qquad
   \ce{N}: \frac{17.27}{14.01} = 1.233 \]

\textbf{Step 3 --- divide by the smallest (1.233):}
\[ \ce{C}: 5.00 \qquad \ce{H}: 7.00 \qquad \ce{N}: 1.00 \]

\textbf{Step 4 --- already whole:} empirical formula \ce{C5H7N},
$M_{\text{emp}} = 5(12.01) + 7(1.008) + 14.01 =
\SI{81.12}{\gram\per\mole}$.
\[ \frac{162.23}{81.12} = 2.00 \qquad\Longrightarrow\qquad
   \textbf{molecular formula } \ce{C10H14N2} \]

Note what the molar mass was \emph{for}: the percentages alone can never
distinguish \ce{C5H7N} from \ce{C10H14N2} --- both have identical
composition. The empirical formula is what analysis gives; the molar mass
picks the molecule.
\end{workedexample}

\begin{workedexample}[Worked example 2: combustion analysis --- a hydrocarbon]
Burning \SI{0.7210}{\gram} of a compound containing only C and H produces
\SI{2.201}{\gram} of \ce{CO2} and \SI{1.081}{\gram} of \ce{H2O}. The molar
mass is \SI{72.15}{\gram\per\mole}. Find both formulas.

\textbf{The idea:} in excess oxygen, every carbon atom ends up in a
\ce{CO2} and every hydrogen atom in an \ce{H2O}. Count them there.

\textbf{Carbon} (one C per \ce{CO2}):
\[ \frac{2.201}{44.01} = \SI{0.05001}{\mole}~\ce{CO2}
   \;\Rightarrow\; \SI{0.05001}{\mole}~\ce{C}
   \;\Rightarrow\; 0.05001 \times 12.01 = \SI{0.6006}{\gram}~\ce{C} \]

\textbf{Hydrogen} (\emph{two} H per \ce{H2O} --- the factor everyone
forgets):
\[ \frac{1.081}{18.02} = \SI{0.05999}{\mole}~\ce{H2O}
   \;\Rightarrow\; \SI{0.1200}{\mole}~\ce{H}
   \;\Rightarrow\; 0.1200 \times 1.008 = \SI{0.1210}{\gram}~\ce{H} \]

\textbf{Mass check:} $0.6006 + 0.1210 = \SI{0.7216}{\gram}$ --- the whole
sample is accounted for, so C and H really are the only elements.

\textbf{Ratio:} $\ce{H}/\ce{C} = 0.1200/0.05001 = 2.40 = 12/5$, so the
empirical formula is \ce{C5H12} with
$M_{\text{emp}} = \SI{72.15}{\gram\per\mole}$ --- equal to the molar mass,
so the \textbf{molecular formula is also \ce{C5H12}} (pentane). The
multiplier can be 1; do not assume it is bigger.
\end{workedexample}

\begin{workedexample}[Worked example 3: combustion with oxygen in the compound]
Burning \SI{2.000}{\gram} of a compound containing C, H, and O gives
\SI{2.868}{\gram} \ce{CO2} and \SI{1.566}{\gram} \ce{H2O}. The molar mass
is \SI{92.09}{\gram\per\mole}. Find both formulas.

\textbf{C and H exactly as before:}
\[ \ce{C}: \frac{2.868}{44.01} = \SI{0.06517}{\mole}
   \to \SI{0.7827}{\gram} \qquad
   \ce{H}: 2 \times \frac{1.566}{18.02} = \SI{0.1738}{\mole}
   \to \SI{0.1752}{\gram} \]

\textbf{Oxygen must come by difference.} The sample's own oxygen atoms end
up in the same \ce{CO2} and \ce{H2O} as the oxygen from the air, so no
measurement can see them separately:
\[ m_{\ce{O}} = 2.000 - 0.7827 - 0.1752 = \SI{1.042}{\gram}
   \;\Rightarrow\; \frac{1.042}{16.00} = \SI{0.06513}{\mole} \]

\textbf{Divide by the smallest (0.06513):}
\[ \ce{C}: 1.00 \qquad \ce{H}: 2.67 \qquad \ce{O}: 1.00 \]

2.67 is $8/3$ --- a value like this must \textbf{not} be rounded. Multiply
all three by 3:
\[ \text{empirical } \ce{C3H8O3}, \quad M_{\text{emp}} = \SI{92.09}{\gram\per\mole}
   = M \;\Rightarrow\; \textbf{molecular } \ce{C3H8O3}~\text{(glycerol)} \]
\end{workedexample}

\begin{apctrap}
The rounding rule, stated once and used forever: ratios such as 1.99 or
3.02 round to whole numbers; ratios sitting near a simple fraction ---
2.25 ($=9/4$), 2.33 ($=7/3$), 2.50, 2.67 ($=8/3$), 2.75 --- are
instructions to \emph{multiply}, by 4, 3, 2, 3, 4 respectively. Rounding
2.67 to 3 produces a wrong formula that no later step can rescue.
\end{apctrap}

\subsection*{\sffamily Check yourself --- \S3.7}

\begin{enumerate}[leftmargin=*,itemsep=0.4em]
  \item A compound is 92.26\% C and 7.74\% H with molar mass
        \SI{78.11}{\gram\per\mole}. Both formulas: \workspace[2.2cm]
        \keyonly{\begin{apcanswer}C: $92.26/12.01 = 7.682$;
        H: $7.74/1.008 = 7.68$. Ratio $1{:}1 \to$ \ce{CH}
        ($M_{\text{emp}} = 13.02$); $78.11/13.02 = 6 \to$
        \ce{C6H6} (benzene).\end{apcanswer}}
  \item Why does combustion analysis measure oxygen by difference rather
        than directly? \answerlines[2]{The compound's oxygen and the
        oxygen supplied for burning both end up in the same \ce{CO2} and
        \ce{H2O}, so they cannot be told apart; only the C and H masses
        are measurable, and O is what remains of the sample mass.}
  \item A mole-ratio calculation gives C 1.00, H 1.33, O 1.00. The
        empirical formula is: \answerblank[3.2cm]{\ce{C3H4O3}}
\end{enumerate}

% =====================================================================
\section*{\sffamily \S3.8 \textbullet\ What a chemical equation says}
\zsec{3.8}

A chemical equation is a statement about \emph{atoms rearranging}:

\[ \ce{CH4(g) + 2O2(g) -> CO2(g) + 2H2O(g)} \]

\begin{center}
\renewcommand{\arraystretch}{1.25}
\begin{tabular}{@{}lll@{}}
\hline
\textbf{Symbol} & \textbf{Means} & \\
\hline
\ce{(s)}, \ce{(l)}, \ce{(g)} & solid, liquid, gas & \\
\ce{(aq)} & dissolved in water & \\
coefficient & \emph{relative} number of molecules or moles &
  never a fixed amount \\
\hline
\end{tabular}
\end{center}

What the equation above claims: one \ce{CH4} molecule and two \ce{O2}
molecules always rearrange into one \ce{CO2} and two \ce{H2O}. The same
sentence read in moles: 1 mol \ce{CH4} reacts with 2 mol \ce{O2}. It does
\emph{not} claim you have 1 mol of anything --- coefficients are ratios,
not inventory.

\textbf{Conserved:} atoms of each element, and therefore total mass.
\textbf{Not conserved:} molecules (3 become 3 here, but 4 become 3 in
ammonia synthesis) and moles. Only atoms are bookkept.

\begin{apctrap}
\textbf{A coefficient multiplies the whole formula; a subscript is part of
the substance's identity.} Balancing is done \emph{only} with
coefficients. ``Fixing'' \ce{H2O} to \ce{H2O2} to balance oxygen does not
balance the equation --- it changes water into hydrogen peroxide and
describes a different reaction altogether.
\end{apctrap}

\subsection*{\sffamily Check yourself --- \S3.8}

\begin{enumerate}[leftmargin=*,itemsep=0.4em]
  \item In \ce{2H2 + O2 -> 2H2O}, four molecules become two. Is mass
        conserved? Explain in one sentence.
        \answerlines[2]{Yes --- every H and O atom on the left appears on
        the right; molecule count is simply not a conserved quantity.}
  \item Write what \ce{2Na(s) + 2H2O(l) -> 2NaOH(aq) + H2(g)} says about
        states. \answerlines[2]{Solid sodium reacts with liquid water to
        give sodium hydroxide dissolved in the water and hydrogen gas.}
\end{enumerate}

% =====================================================================
\section*{\sffamily \S3.9 \textbullet\ Balancing chemical equations}
\zsec{3.9}

Balancing is done \textbf{by inspection}, but inspection with a strategy:

\begin{enumerate}[leftmargin=*,itemsep=0.2em]
  \item Start from the \emph{most complicated} molecule.
  \item Balance elements that appear in only one place on each side first.
  \item Leave lone elements (\ce{O2}, \ce{H2}, metals) for \emph{last} ---
        their coefficient can be set freely without unbalancing anything
        else.
  \item Count every atom on both sides at the end. Always.
\end{enumerate}

\begin{workedexample}[Worked example 4: ethanol combustion, step by step]
\[ \ce{C2H5OH(l) + O2(g) -> CO2(g) + H2O(g)} \qquad \text{(unbalanced)} \]

\textbf{Start with ethanol} (most complex, sets C and H):
carbon first --- 2 C on the left forces
\[ \ce{C2H5OH + O2 -> 2CO2 + H2O} \]
hydrogen next --- 6 H on the left ($5+1$) forces
\[ \ce{C2H5OH + O2 -> 2CO2 + 3H2O} \]

\textbf{Oxygen last}, because \ce{O2} stands alone. The right side now has
$2(2) + 3(1) = 7$ O; the left has 1 in the ethanol, so \ce{O2} must supply
6 atoms $= 3$ molecules:
\[ \boxed{\;\ce{C2H5OH(l) + 3O2(g) -> 2CO2(g) + 3H2O(g)}\;} \]

\textbf{Final count:} C $2 = 2$; H $6 = 6$; O $1 + 6 = 4 + 3$. Balanced.
\end{workedexample}

\begin{workedexample}[Worked example 5: when a fraction appears --- hexane]
\[ \ce{C6H14(l) + O2(g) -> CO2(g) + H2O(g)} \]
C and H from the hexane: $6\,\ce{CO2}$ and $7\,\ce{H2O}$. The right side
then carries $12 + 7 = 19$ oxygen atoms --- an \emph{odd} number, which a
whole coefficient on \ce{O2} cannot supply. Balance with a fraction first:
\[ \ce{C6H14 + \tfrac{19}{2}O2 -> 6CO2 + 7H2O} \]
then clear it by doubling \textbf{every} coefficient:
\[ \boxed{\;\ce{2C6H14(l) + 19O2(g) -> 12CO2(g) + 14H2O(g)}\;} \]
Count: C $12 = 12$; H $28 = 28$; O $38 = 24 + 14$. The doubling step must
touch all four coefficients --- doubling only the \ce{O2} is the standard
error.
\end{workedexample}

\subsection*{\sffamily Check yourself --- \S3.9}

Balance each (smallest whole coefficients):

\begin{enumerate}[leftmargin=*,itemsep=0.4em]
  \item \ce{Fe(s) + O2(g) -> Fe2O3(s)}
        \answerblank[6cm]{\ce{4Fe + 3O2 -> 2Fe2O3}}
  \item \ce{C3H8(g) + O2(g) -> CO2(g) + H2O(g)}
        \answerblank[7.2cm]{\ce{C3H8 + 5O2 -> 3CO2 + 4H2O}}
  \item \ce{Al(s) + HCl(aq) -> AlCl3(aq) + H2(g)}
        \answerblank[7.2cm]{\ce{2Al + 6HCl -> 2AlCl3 + 3H2}}
  \item \ce{C4H10(g) + O2(g) -> CO2(g) + H2O(g)}
        \answerblank[8cm]{\ce{2C4H10 + 13O2 -> 8CO2 + 10H2O}}
\end{enumerate}

% =====================================================================
\section*{\sffamily \S3.10 \textbullet\ Stoichiometric calculations}
\zsec{3.10}

The balanced equation is a \emph{mole-to-mole} conversion factor and
nothing more. Every mass problem is the same three moves:

\begin{center}
\fbox{\parbox{0.9\linewidth}{\centering
grams of A $\xrightarrow{\;\div M_A\;}$ moles of A
$\xrightarrow{\;\text{coefficients}\;}$ moles of B
$\xrightarrow{\;\times M_B\;}$ grams of B}}
\end{center}

\begin{workedexample}[Worked example 6: burning propane end to end]
A grill burns \SI{25.0}{\gram} of propane:
$\ce{C3H8(g) + 5O2(g) -> 3CO2(g) + 4H2O(g)}$.
What mass of \ce{O2} is consumed, and what mass of \ce{CO2} is produced?

\textbf{To moles once, at the start:}
\[ \frac{25.0}{44.09} = \SI{0.567}{\mole}~\ce{C3H8} \]

\textbf{Mole ratios, straight off the coefficients:}
\[ \ce{O2}: 0.567 \times \frac{5}{1} = \SI{2.835}{\mole} \qquad
   \ce{CO2}: 0.567 \times \frac{3}{1} = \SI{1.701}{\mole} \]

\textbf{Back to grams at the end:}
\[ m_{\ce{O2}} = 2.835 \times 32.00 = \mathbf{90.7~g} \qquad
   m_{\ce{CO2}} = 1.701 \times 44.01 = \mathbf{74.9~g} \]

\textbf{Sanity check by conservation of mass:} reactants
$25.0 + 90.7 = 115.7$~g; products $74.9 + $ water. Water:
$0.567 \times 4 = 2.268$~mol $\times\, 18.02 = 40.9$~g, and
$74.9 + 40.9 = 115.8$~g. The books balance (0.1 g is rounding). Thirty
seconds of checking catches almost every slip in this chapter.
\end{workedexample}

\begin{workedexample}[Worked example 7: working backward --- iron from ore]
Iron is produced in a blast furnace by
$\ce{Fe2O3(s) + 3CO(g) -> 2Fe(l) + 3CO2(g)}$.
What mass of \ce{Fe2O3} is required for \SI{1.00}{\kilo\gram} of iron?

\[ \frac{1000}{55.85} = \SI{17.91}{\mole}~\ce{Fe}
   \;\xrightarrow{\times 1/2}\; \SI{8.953}{\mole}~\ce{Fe2O3}
   \;\xrightarrow{\times 159.70}\; \mathbf{1.43\times10^{3}~g}
   = \mathbf{1.43~kg} \]

The ratio ran ``backwards'' (product to reactant) with no change in
method --- the equation converts in either direction. Note the ratio is
$\tfrac{1~\ce{Fe2O3}}{2~\ce{Fe}}$, taken \emph{from the coefficients},
never from the formulas' subscripts.
\end{workedexample}

\subsection*{\sffamily Check yourself --- \S3.10}

\begin{enumerate}[leftmargin=*,itemsep=0.4em]
  \item Using worked example 6's equation, what mass of water forms when
        \SI{10.0}{\gram} of propane burns? \workspace[2.2cm]
        \keyonly{\begin{apcanswer}$10.0/44.09 = 0.227$~mol;
        $\times 4 = 0.907$~mol \ce{H2O}; $\times 18.02 =
        \mathbf{16.3}$~g\end{apcanswer}}
  \item In \ce{2Al + 3Cl2 -> 2AlCl3}, how many moles of \ce{Cl2} react
        with \SI{0.40}{\mole} Al? \answerblank[2.6cm]{0.60 mol}
  \item A student converts \SI{25.0}{\gram} \ce{C3H8} to grams of
        \ce{CO2} by $25.0 \times \tfrac{3}{1}$. Name the two omitted
        steps. \answerlines[2]{Grams were never converted to moles, and
        the result was never converted back --- the mole ratio was applied
        to masses. Grams never talk to grams directly.}
\end{enumerate}

% =====================================================================
\section*{\sffamily \S3.11 \textbullet\ Limiting reactant and percent
yield}
\zsec{3.11}

Real mixtures are not delivered in perfect stoichiometric ratio. Whichever
reactant runs out first --- the \term{limiting reactant} --- fixes how much
product can form; everything else is \term{in excess} and some of it is
left over.

\textbf{To find the limiting reactant}, convert every reactant to moles,
then pick either test:

\begin{enumerate}[leftmargin=*,itemsep=0.2em]
  \item \textbf{Required-vs-available.} Use the mole ratio to compute how
        much of B reactant A would need; compare with what is present.
  \item \textbf{Smallest-product.} Compute the product each reactant could
        make alone; the smaller answer is the truth, and the reactant that
        produced it is limiting.
\end{enumerate}

Both always agree. Use whichever reads more naturally, but \emph{never}
compare raw masses or raw moles of the reactants --- the comparison only
means anything after the coefficients are involved.

\begin{workedexample}[Worked example 8: ammonia synthesis, complete]
\SI{50.0}{\gram} of \ce{N2} is mixed with \SI{15.0}{\gram} of \ce{H2}:
$\ce{N2(g) + 3H2(g) -> 2NH3(g)}$. Find the limiting reactant, the
theoretical yield of \ce{NH3}, and the mass of the excess reactant left
over.

\textbf{Everything to moles first:}
\[ \ce{N2}: \frac{50.0}{28.02} = \SI{1.784}{\mole} \qquad
   \ce{H2}: \frac{15.0}{2.016} = \SI{7.44}{\mole} \]

\textbf{Test 1 (required vs available):} the \ce{N2} present would need
$1.784 \times 3 = \SI{5.35}{\mole}$ of \ce{H2}. Available: 7.44 mol ---
more than needed, so \ce{H2} is in excess and \textbf{\ce{N2} is
limiting}. (Notice the raw moles pointed the other way: there is
\emph{more} \ce{H2} by moles, yet \ce{N2} still limits. The 3:1 ratio is
what decides.)

\textbf{Theoretical yield} --- from the limiting reactant only:
\[ 1.784 \times \frac{2~\ce{NH3}}{1~\ce{N2}} = \SI{3.569}{\mole}
   \;\Rightarrow\; 3.569 \times 17.03 = \mathbf{60.8~g~\ce{NH3}} \]

\textbf{Leftover excess:} \ce{H2} consumed $= 5.35$~mol, so
\[ 7.44 - 5.35 = \SI{2.09}{\mole} \;\Rightarrow\;
   2.09 \times 2.016 = \mathbf{4.21~g~\ce{H2}~left} \]

\textbf{Mass check:} in $50.0 + 15.0 = 65.0$~g; out
$60.8 + 4.21 = 65.0$~g. Balanced books again.
\end{workedexample}

\begin{workedexample}[Worked example 9: percent yield]
The synthesis in worked example 8 is run and \SI{48.0}{\gram} of \ce{NH3}
is actually collected.
\[ \%~\text{yield} = \frac{\text{actual}}{\text{theoretical}} \times 100
   = \frac{48.0}{60.8} \times 100 = \mathbf{78.9\%} \]
The \term{theoretical yield} is a calculation; the \term{actual yield} is
a measurement, and it is smaller for ordinary reasons --- side reactions,
incomplete reaction, product lost in handling. A percent yield over 100\%
is not good news: it means the product is wet or impure, or the
bookkeeping is wrong.
\end{workedexample}

\begin{apctrap}
Three habits that prevent nearly every limiting-reactant error:
theoretical yield is computed from the \textbf{limiting} reactant only,
never from both; the leftover is
$(\text{present}) - (\text{consumed})$, where consumed comes through the
mole ratio; and the limiting reactant is \textbf{not} the one with fewer
grams or fewer moles --- \SI{50.0}{\gram} of \ce{N2} limited against
\SI{15.0}{\gram} of \ce{H2} above.
\end{apctrap}

\subsection*{\sffamily Check yourself --- \S3.11}

\begin{enumerate}[leftmargin=*,itemsep=0.4em]
  \item \SI{5.40}{\gram} Al and \SI{8.10}{\gram} \ce{Cl2} react:
        \ce{2Al + 3Cl2 -> 2AlCl3}. Identify the limiting reactant and the
        theoretical yield of \ce{AlCl3}. \workspace[3cm]
        \keyonly{\begin{apcanswer}Al: $5.40/26.98 = 0.2002$~mol;
        \ce{Cl2}: $8.10/70.90 = 0.1142$~mol. The Al would need
        $0.2002 \times \tfrac{3}{2} = 0.300$~mol \ce{Cl2}; only 0.1142 is
        present, so \textbf{\ce{Cl2} is limiting}.
        \ce{AlCl3} $= 0.1142 \times \tfrac{2}{3} = 0.0762$~mol
        $\times\, 133.33 = \mathbf{10.2}$~g.\end{apcanswer}}
  \item For that mixture, the mass of aluminum left over:
        \workspace[2.2cm]
        \keyonly{\begin{apcanswer}Al consumed
        $= 0.1142 \times \tfrac{2}{3} = 0.0762$~mol; left
        $= 0.2002 - 0.0762 = 0.1240$~mol $\times\, 26.98 =
        \mathbf{3.35}$~g.\end{apcanswer}}
  \item If \SI{9.20}{\gram} of \ce{AlCl3} is collected, the percent
        yield (divide by the \emph{unrounded} theoretical yield):
        \answerblank[2.6cm]{90.6\%}
  \item True or false, with a reason: the reactant present in the
        smaller mass is the limiting reactant.
        \answerlines[2]{False --- limitation is decided by moles compared
        through the coefficients. In worked example 8 the smaller mass
        (\ce{H2}) was the excess reactant.}
\end{enumerate}

\begin{apcnote}[The whole span in four sentences]
A formula comes from mole ratios of elements (\S3.7). A reaction is
written as a balanced equation, using coefficients only (\S3.8--3.9). The
equation's coefficients convert moles of anything into moles of anything
else, with grams only at the edges (\S3.10). When amounts of two reactants
are given, the mole ratio decides which one limits, and everything ---
yield and leftovers --- follows from that one reactant (\S3.11).
\end{apcnote}

\end{document}
